% Source: https://artofproblemsolving.com/community/c3h2337421_what_the_answer_to_this_problem_
% Stable identity: aops-topic-2337421-q01
Following the order of operations, first $(1-6)=-5$ and $2\cdot3=6$; then add and subtract from left to right:
\[-2-3+(-5)+6=-5+(-5)+6=-10+6=\boxed{-4}.\]
So Way 1 is correct.

Way 2 goes wrong twice, both times by letting a minus sign act on more than the term right after it.
\begin{itemize}
\item In $-2-(+3)+(-5)$, the minus sign subtracts only $3$: $-(+3)+(-5)=-3-5=-8$. Way 2 replaces it by $-\bigl((+3)+(-5)\bigr)=-(-2)=+2$, as if the minus sign applied to the sum.
\item In $-2-(-2)+6$ the same mistake is made again: Way 2 writes $-(-2)+6$ as $-(+4)$, i.e. $-\bigl((-2)+6\bigr)$, whereas $-(-2)+6=2+6=8$.
\end{itemize}
(Subtraction is not associative: $a-b+c=(a-b)+c$, not $a-(b+c)$.)
